\section{Conclusion}
\label{sec:conclusion}

This study examined how replay-system choices interact under a fixed-total
exemplar budget for CIFAR-100 class-incremental learning.  Holding the protocol
and resource budget fixed made it possible to separate exemplar selection,
distillation, classifier head, and evaluation readout.  The matched ablations
show that nearest-mean exemplar evaluation and herding are the largest measured
contributors to the reference configuration's accuracy--forgetting trade-off.
Distillation primarily improves retention, whereas the memory sweep has a
larger effect than the tested changes in retrieval budget.

For the measured setting, a defensible replay design is to allocate memory
across all observed classes, rebuild representative exemplars in the current
feature space, train with distillation, and evaluate with the corresponding NME
prototypes.  This is a protocol-specific recommendation, not a general ranking
of class-incremental methods.  The next empirical step is to test the same
component-level analysis across class orders, datasets, architectures, and
larger memory ranges.
